% B01-S001 | APPM5720Notes.tex baris 1263-1321
% segment-id: d90.becker.98ed693.b01.seg0001
\chapter{Dualitas Lagrange, kondisi Slater, dan kondisi KKT}
\label{becker:chapter:lagrange-kkt}

Bab ini melengkapi tulang punggung Habring dengan satu jalur terpadu dari
Lagrangian ke kondisi Slater dan kondisi Karush-Kuhn-Tucker (KKT). Notasi
dibuat konsisten untuk masalah minimisasi. Bagian sumber yang khusus membahas
dualitas program linear tidak disertakan karena materi itu berada dalam batas
kurikulum O018.

\noindent\textbf{Kredit catatan donor.}\par
Catatan donor menyatakan bahwa bagian
dualitas mengikuti Bab~5 Boyd dan Vandenberghe (B\&V), merujuk ke
\S5.3 untuk ilustrasi, dan mengaitkan interpretasi perpotongan maksimum atau
titik bersama minimum dengan Bertsekas. Kredit turunan ini dipertahankan;
edisi ini tidak mengimpor isi di luar rentang donor yang dibekukan.

\section{Masalah primal dan Lagrangian}

Pertimbangkan masalah
\begin{equation}
  \label{becker:eq:primal}
  \begin{aligned}
    p^* = \inf_{x\in D}\quad & f_0(x)\\
    \text{dengan syarat}\quad & f_i(x)\leq 0, && i=1,\ldots,m,\\
    & h_j(x)=0, && j=1,\ldots,p,
  \end{aligned}
\end{equation}
dengan $D$ memuat domain bersama fungsi-fungsi tersebut. Kita memakai
$p^*=+\infty$ jika masalah tidak layak dan $p^*=-\infty$ jika nilainya tidak
terbatas ke bawah.

\begin{defn}[Lagrangian]
  Lagrangian masalah~\eqref{becker:eq:primal} adalah
  \begin{equation}
    \label{becker:eq:lagrangian}
    L(x,\lambda,\nu)
    = f_0(x)+\sum_{i=1}^{m}\lambda_i f_i(x)
      +\sum_{j=1}^{p}\nu_j h_j(x).
  \end{equation}
  Vektor $\lambda\in\R^m$ dan $\nu\in\R^p$ disebut variabel dual atau
  pengali Lagrange. Untuk kendala pertidaksamaan kita akan mensyaratkan
  $\lambda\geq0$ komponen demi komponen.
\end{defn}

Lagrangian bergantung pada cara masalah ditulis. Misalnya, untuk
$\inf_x f(x)+g(x)$ kita dapat memperkenalkan $z$ dan kendala $x=z$, lalu
membentuk Lagrangian terhadap kendala kesamaan itu. Pemisahan variabel seperti
ini sering membuka struktur yang tidak tampak pada rumusan asal.

\begin{defn}[Fungsi dual]
  Fungsi dual Lagrange didefinisikan oleh
  \begin{equation}
    \label{becker:eq:dual-function}
    g(\lambda,\nu)=\inf_{x\in D} L(x,\lambda,\nu).
  \end{equation}
  Fungsi $g$ cekung terhadap $(\lambda,\nu)$, bahkan bila masalah primal
  tidak konveks, sebab $g$ adalah infimum titik demi titik dari fungsi-fungsi
  afin dalam variabel dual.
\end{defn}

Sebagai kaidah umum yang terkait, minimisasi parsial dari fungsi konveks
bersama menghasilkan fungsi konveks pada variabel yang tersisa, sedangkan
supremum titik demi titik dari fungsi-fungsi konveks tetap konveks.

\begin{defn}[Masalah dual]
  Masalah dual Lagrange adalah
  \begin{equation}
    \label{becker:eq:dual-problem}
    d^*=\sup_{\lambda\geq0,\,\nu\in\R^p} g(\lambda,\nu).
  \end{equation}
  Ini merupakan masalah maksimisasi cekung, atau secara ekuivalen masalah
  minimisasi konveks untuk $-g$.
\end{defn}

\begin{theorem}[Dualitas lemah]
  Untuk masalah~\eqref{becker:eq:primal} selalu berlaku
  \begin{equation}
    \label{becker:eq:weak-duality}
    d^*\leq p^*.
  \end{equation}
\end{theorem}
\begin{proof}
  Ambil titik layak $\tilde x$. Untuk setiap $\lambda\geq0$ dan setiap
  $\nu$,
  \begin{multline}
    g(\lambda,\nu)\leq L(\tilde x,\lambda,\nu)
    =f_0(\tilde x)
      +\underbrace{\sum_{i=1}^{m}\lambda_i f_i(\tilde x)}_{\leq0}
      +\underbrace{\sum_{j=1}^{p}\nu_j h_j(\tilde x)}_{=0}\\
    \leq f_0(\tilde x).
  \end{multline}
  Mengambil infimum terhadap semua titik primal yang layak memberi
  $g(\lambda,\nu)\leq p^*$. Mengambil supremum terhadap semua titik dual yang
  layak membuktikan~\eqref{becker:eq:weak-duality}.
\end{proof}

Jika $d^*=p^*$, kita mengatakan bahwa \emph{dualitas kuat} berlaku. Selain
memberi sertifikat batas bawah, masalah dual dapat memindahkan operator linear,
menyingkap kehalusan, atau menghasilkan rumusan numerik yang lebih sesuai.
Untuk masalah nonkonveks sekalipun, nilai dual tetap menyediakan batas bawah
yang sah, walaupun kesenjangan dualitas mungkin positif.

% B01-S002 | APPM5720Notes.tex baris 1398-1405 dan 1414-1470
% segment-id: d90.becker.98ed693.b01.seg0002
\section{Kondisi Slater dan dualitas kuat}

Sekarang andaikan $f_0,f_1,\ldots,f_m$ konveks dan kendala kesamaan berbentuk
$Ax=b$. Tuliskan $I_{\mathrm{na}}$ bagi indeks kendala pertidaksamaan yang tidak
afin.

\begin{defn}[Kondisi Slater]
  Kondisi Slater dipenuhi jika terdapat $\bar x$ dalam interior relatif domain
  bersama sedemikian sehingga
  \begin{equation}
    A\bar x=b,\qquad
    f_i(\bar x)<0\quad(i\in I_{\mathrm{na}}),\qquad
    f_i(\bar x)\leq0\quad(i\notin I_{\mathrm{na}}).
  \end{equation}
  Jadi hanya kendala pertidaksamaan nonafin yang harus dipenuhi secara ketat.
\end{defn}

\begin{theorem}[Slater]
  Jika masalah primal konveks, kondisi Slater dipenuhi, dan $p^*$ berhingga,
  maka $d^*=p^*$ dan suatu titik dual optimum $(\lambda^*,\nu^*)$ tercapai.
\end{theorem}

Kondisi Slater bersifat cukup, bukan perlu. Contoh semidefinit berikut
menunjukkan salah satu gejala degenerasi. Pertimbangkan
\begin{equation}
  \inf_{X=X^\top\succeq0}
  \left\langle
    \begin{pmatrix}1&0\\0&0\end{pmatrix},X
  \right\rangle
  \quad\text{dengan syarat}\quad
  \left\langle
    \begin{pmatrix}0&1\\1&0\end{pmatrix},X
  \right\rangle=2.
\end{equation}
Setiap titik layak berbentuk
$X=\begin{psmallmatrix}a&1\\1&b\end{psmallmatrix}\succeq0$, sehingga
$a,b\geq0$ dan $ab\geq1$. Barisan
$X_\varepsilon=\begin{psmallmatrix}\varepsilon&1\\1&1/\varepsilon\end{psmallmatrix}$
memberi nilai objektif yang menuju nol ketika $\varepsilon\downarrow0$.
Jadi $p^*=0$, tetapi infimum primal tidak tercapai. Dualnya tetap mempunyai
solusi optimum bernilai nol; contoh ini memperingatkan bahwa ketercapaian
primal tidak mengikuti hanya dari kesamaan nilai primal dan dual.

Berikut gagasan geometri di balik teorema Slater. Definisikan himpunan nilai
terganggu
\begin{equation}
  \mathcal A=\left\{(u,v,t):\ \exists x\in D,\quad
  f_i(x)\leq u_i,\ h(x)=v,\ f_0(x)\leq t\right\}
\end{equation}
dan sinar terbuka
\begin{equation}
  \mathcal B=\{(0,0,s):s<p^*\}.
\end{equation}
Kekonveksan fungsi-fungsi kendala membuat $\mathcal A$ konveks, dan
$\mathcal A\cap\mathcal B=\varnothing$. Teorema pemisahan memberi normal
$(\lambda^*,\nu^*,\mu^*)$. Monotonisitas $\mathcal A$ terhadap $u$ dan $t$
memberi $\lambda^*\geq0$ dan $\mu^*\geq0$. Kondisi Slater menyingkirkan
$\mu^*=0$, yaitu hiperbidang penyangga vertikal, sehingga $\mu^*>0$ dan normal
dapat dinormalkan menjadi
$(\lambda^*,\nu^*,1)$. Ketaksamaan pemisahan kemudian menghasilkan
\begin{equation}
  \inf_{x\in D}L(x,\lambda^*,\nu^*)\geq p^*.
\end{equation}
Dualitas lemah memberi ketaksamaan sebaliknya; karena itu $d^*=p^*$ dan
pengali dual optimum tercapai. Interpretasi yang sama dapat dinyatakan sebagai
perpotongan maksimum atau titik bersama minimum dalam geometri perturbasi.

% B01-S003 | APPM5720Notes.tex baris 1472-1499
% segment-id: d90.becker.98ed693.b01.seg0003
\section{Interpretasi titik pelana}

Untuk Lagrangian~\eqref{becker:eq:lagrangian}, masalah primal dan dual dapat
ditulis sebagai
\begin{equation}
  \begin{aligned}
    p^*&=\inf_x\sup_{\lambda\geq0,\nu}L(x,\lambda,\nu),\\
    d^*&=\sup_{\lambda\geq0,\nu}\inf_xL(x,\lambda,\nu).
  \end{aligned}
\end{equation}
Ketaksamaan minimaks
\begin{equation}
  \sup_{\lambda\geq0,\nu}\inf_xL(x,\lambda,\nu)
  \leq
  \inf_x\sup_{\lambda\geq0,\nu}L(x,\lambda,\nu)
\end{equation}
adalah bentuk titik-pelana dari dualitas lemah.

\begin{defn}[Titik pelana]
  Tripel $(x^*,\lambda^*,\nu^*)$ adalah titik pelana Lagrangian jika
  $\lambda^*\geq0$ dan
  \begin{equation}
    L(x^*,\lambda,\nu)
    \leq L(x^*,\lambda^*,\nu^*)
    \leq L(x,\lambda^*,\nu^*)
  \end{equation}
  untuk setiap $x\in D$, $\lambda\geq0$, dan $\nu\in\R^p$.
\end{defn}

Sebagai contoh, pertimbangkan
\begin{equation}
  \min_x\norm{x}_1
  \quad\text{dengan syarat}\quad
  \norm{Ax-b}_2^2\leq\varepsilon^2.
\end{equation}
Lagrangiannya adalah
\begin{equation}
  L(x,\lambda)=\norm{x}_1
  +\lambda\bigl(\norm{Ax-b}_2^2-\varepsilon^2\bigr),\qquad\lambda\geq0.
\end{equation}
Jika pengali optimum memenuhi $\lambda^*>0$, minimisasi Lagrangian terhadap
$x$ ekuivalen, setelah dibagi dengan $2\lambda^*$, dengan
\begin{equation}
  \min_x \frac12\norm{Ax-b}_2^2
  +\frac{1}{2\lambda^*}\norm{x}_1.
\end{equation}
Jadi rumusan berkendala dan rumusan berpenalti berhubungan melalui pengali
optimum. Hubungan ini tidak boleh diperoleh dengan membagi oleh $\lambda^*$
jika pengalinya nol.

% B01-S004 | APPM5720Notes.tex baris 1652-1726
% segment-id: d90.becker.98ed693.b01.seg0004
\section{Kondisi Karush-Kuhn-Tucker}

Andaikan fungsi-fungsi pada masalah~\eqref{becker:eq:primal} dapat
didiferensialkan. Tripel $(x^*,\lambda^*,\nu^*)$ memenuhi kondisi KKT jika
\begin{enumerate}
  \item \textbf{stasioneritas:}
  $\nabla_xL(x^*,\lambda^*,\nu^*)=0$;
  pada kasus konveks tak mulus, bentuk yang sesuai adalah
  $0\in\partial_xL(x^*,\lambda^*,\nu^*)$;
  \item \textbf{kelayakan primal:}
  $f_i(x^*)\leq0$ dan $h_j(x^*)=0$;
  \item \textbf{kelayakan dual:} $\lambda^*\geq0$;
  \item \textbf{kekomplementeran:}
  $\lambda_i^*f_i(x^*)=0$ untuk setiap $i=1,\ldots,m$.
\end{enumerate}

\begin{theorem}[KKT perlu dari pasangan primal-dual optimum]
  Andaikan optimum primal dan dual tercapai, dualitas kuat berlaku, dan
  $x\mapsto L(x,\lambda^*,\nu^*)$ memiliki syarat optimalitas Fermat yang
  berlaku pada $x^*$. Maka pasangan optimum memenuhi kondisi KKT. Pernyataan
  ini tetap dapat dipakai pada masalah nonkonveks bila stasioneritas Lagrangian
  benar-benar mengikuti dari minimisasi globalnya dan seluruh syarat domain
  dipenuhi.
\end{theorem}

\begin{theorem}[KKT cukup untuk masalah konveks]
  Jika $f_0,f_1,\ldots,f_m$ konveks, $h_j$ afin, dan suatu tripel memenuhi
  kondisi KKT, maka $x^*$ optimum primal, $(\lambda^*,\nu^*)$ optimum dual,
  dan $p^*=d^*$.
\end{theorem}
\begin{proof}
  Stasioneritas dan kekonveksan
  $x\mapsto L(x,\lambda^*,\nu^*)$ membuat $x^*$ peminimum global Lagrangian.
  Karena itu
  \begin{equation}
    \begin{aligned}
      d^*&\geq g(\lambda^*,\nu^*)
        =\inf_xL(x,\lambda^*,\nu^*)\\
      &=L(x^*,\lambda^*,\nu^*)
        =f_0(x^*)
         +\sum_i\lambda_i^*f_i(x^*)
         +\sum_j\nu_j^*h_j(x^*)\\
      &=f_0(x^*)\geq p^*.
    \end{aligned}
  \end{equation}
  Dualitas lemah memberi $d^*\leq p^*$, sehingga semua ketaksamaan adalah
  kesamaan.
\end{proof}

\begin{theorem}[KKT perlu di bawah kondisi Slater]
  Jika masalah primal konveks memenuhi kondisi Slater dan suatu solusi primal
  $x^*$ ada dengan nilai berhingga, maka terdapat pengali
  $(\lambda^*,\nu^*)$ sedemikian sehingga tripel tersebut memenuhi kondisi
  KKT.
\end{theorem}

\subsection{Contoh: proyeksi pada bola $\ell^1$}

Proyeksi $y\in\R^n$ pada bola
$B_1(\tau)=\{x:\norm{x}_1\leq\tau\}$, dengan $\tau>0$, menyelesaikan
\begin{equation}
  \min_x \frac12\norm{x-y}_2^2
  \quad\text{dengan syarat}\quad \norm{x}_1\leq\tau.
\end{equation}
Kondisi KKT memberi
\begin{equation}
  x_i=S_\lambda(y_i)
  =\operatorname{sign}(y_i)\max\{|y_i|-\lambda,0\},
  \qquad \lambda\geq0,
\end{equation}
bersama $\norm{x}_1\leq\tau$ dan
$\lambda(\norm{x}_1-\tau)=0$. Jika $\norm{y}_1\leq\tau$, maka
$\lambda=0$ dan $x=y$. Jika $\norm{y}_1>\tau$, maka $\lambda>0$ adalah akar
unik dari
\begin{equation}
  \sum_{i=1}^{n}\max\{|y_i|-\lambda,0\}=\tau.
\end{equation}
Ruas kiri kontinu, menurun, dan linear sepotong-sepotong dengan titik patah
$|y_i|$; akar dapat dicari dengan biseksi atau, lebih efisien, dengan
mengurutkan titik patah. Identitas Moreau menjelaskan hubungan yang tepat:
\begin{equation}
  \prox_{\tau\norm{\emptyarg}_\infty}(y)
  =y-\operatorname{proj}_{B_1(\tau)}(y).
\end{equation}

Kekomplementeran juga mengikuti dari kesamaan primal-dual. Untuk titik pelana
optimum $(x^*,\lambda^*,\nu^*)$ berlaku
\begin{equation}
  f_0(x^*)
  =L(x^*,\lambda^*,\nu^*)
  =f_0(x^*)+\sum_i\lambda_i^*f_i(x^*),
\end{equation}
karena kendala kesamaan lenyap. Setiap suku
$\lambda_i^*f_i(x^*)\leq0$, sehingga jumlah nol memaksa
$\lambda_i^*f_i(x^*)=0$ untuk setiap $i$.

% B01-S005 | APPM5720Notes.tex baris 1731-1743
% segment-id: d90.becker.98ed693.b01.seg0005
\subsection{Contoh: masalah kuadratik dengan kendala kesamaan}

Pertimbangkan
\begin{equation}
  \min_x \frac12 x^\top P x+q^\top x+r
  \quad\text{dengan syarat}\quad Ax=b,
  \qquad P\succeq0.
\end{equation}
Kondisi KKT adalah
\begin{equation}
  Px^*+q+A^\top\nu^*=0,
  \qquad Ax^*=b,
\end{equation}
atau sistem titik pelana
\begin{equation}
  \begin{pmatrix}P&A^\top\\A&0\end{pmatrix}
  \begin{pmatrix}x^*\\\nu^*\end{pmatrix}
  =\begin{pmatrix}-q\\b\end{pmatrix}.
\end{equation}
Setiap solusi sistem ini adalah solusi optimum masalah konveks tersebut.
Solusi tunggal hanya dijamin jika matriks KKT nonsingular; satu syarat cukup
yang umum adalah $A$ berperingkat baris penuh dan $P$ definit positif pada
$\ker(A)$.
